\vspace{-0.1in}
\section{Related Work}
\vspace{-0.1in}

\textbf{Learning Rate Schedules.} Learning rate schedules are crucial in deep learning, with previous studies exploring various options. \citet{smith2017cyclical} was the first to propose a cyclic triangular learning rate schedule that interleaves decreasing and increasing learning rates. \citet{loshchilov2017sgdr} extended the idea to a cyclic cosine learning rate schedule. \citet{he2015deep} introduced the notion of warmup, which gradually increases the learning rate in the earlier training phase. \citet{goyal2018accurate, hoffer2018train, you2020large} concluded that the learning rate should scale linearly with the batch size, which is further theoretically examined in~\citet{smith2020generalization, li2021validity, malladi2023sdes}. 
\citet{you2019does} performed an analysis on why learning rate schedules are helpful and suspected that the large learning rate at the beginning phase is mostly useful for avoiding memorization of noisy data, which is consistent with our analysis in~\Cref{sec:data}.

In the LLM era, works including~\citet{hoffmann2022training, deepseekai2024deepseek, hu2024minicpm} examined how to choose learning rate schedules for pretraining. In particular, \citet{hu2024minicpm} introduced a learning rate schedule called Warmup-Stable-Decay ($\wsd$) that remains constant for the majority of the runs before decaying in language model pretraining, which were studied independently in~\citet{zhai2022scaling, ibrahim2024simple, hagele2024scalinglawscomputeoptimaltraining}. \citet{raffel2023exploring, ibrahim2024simple} explored another possibility of using an inverse square root schedule to pretrain the language models. \citet{defazio2023whenmuchadaptivelearning} proposes to use linear decay for the entire training run. \citet{defazio2024roadscheduled} shows that with appropriate iterate averaging, a constant learning rate schedule can reach better performance than the cosine learning rate schedule. \citet{rae2022scaling, gupta2023continual, hu2024minicpm, ibrahim2024simple} examined how to choose a learning rate schedule in a continual learning setting and verified that rewarming-up cosine learning rate brings performance drops that are costly to recover. A common belief is that the performance drop is due to the sudden increase in learning rate during rewarming-up. However, our work shows that increasing the learning rate after a short decay in $\wsd$ does not cause a similar performance drop as seen with the cosine learning rate, challenging the previous hypothesis. Instead, we suggest that the performance loss associated with rewarming-up cosine learning rate is due to the implicit damage it causes to the model, making it unsuitable for continual training. On the contrast, $\wsd$ avoids such damage by maintaining a high learning rate during the stable phase, hence the sudden increase in learning rate does not lead to performance drops in continual training.  


\textbf{Continual Learning.}  Continual learning, the process of updating the model with newly collected data, can improve the models' knowledge and capability.  Previous continual learning research \citep{aljundi2019online, veniat2021efficient, cossu2022continual, 51163, harun2023siesta, mehta2023empirical} assumed  significant domain shift and aimed to avoid forgetting old knowledge while learning new knowledge. 
Recent works including \citet{hernandez2021scaling, lesort2023challenging} suggested that optimizers including SGD and Adam have a knowledge accumulation effect and the effect of catastrophic forgetting may be less significant than expected, especially when replay is applied. Our work mainly focuses on continual pre-training without necessarily a strong domain shift and hence does not touch upon the effect of covariance shift. Continual learning is also extensively employed in large language models such as LLaMA to extend their capabilities, such as handling longer contexts (e.g., see \citet{tworkowski2023focused, peng2023yarnefficientcontextwindow, chen2023extendingcontextwindowlarge, dubey2024llama3herdmodels} and references therein) or dealing with new languages and domains (e.g., see \citet{azerbayev2024llemmaopenlanguagemodel, rozière2024codellamaopenfoundation, cui2024efficienteffectivetextencoding} and references therein).



\textbf{Theoretical Understanding on Loss Landscape.} A long line of research aims to better understand the loss landscape in deep learning (e.g., see \citet{freeman2017topologygeometryhalfrectifiednetwork, garipov2018losssurfacesmodeconnectivity, li2020explaining} and references therein). We will highlight several phenomena that are related to our findings. 

(1) Ill-conditioned directional sharpness and heavy-tailed noise: \citet{zhang2020gradient, NEURIPS2020_b05b57f6} examined the gradient noise in language modeling and observed that the noise is heavy-tailed in multiple dimensions. \citet{pan2023understanding, liu2024sophia} showed that the loss has vastly different curvatures in different dimensions. \citet{pan2022eigencurveoptimallearningrate} analyzes optimizing a quadratic function with skewed curvature theoretically.
Our river valley landscape is consistent with these findings. 

(2) Benefit of large learning rates: Large learning rates have a provable regularizing effect in finding flatter minima~\citep{kong2020stochasticity, wang2022large}, and flatter minima typically have a better generalization effect, even in the pretraining setting~\citep{jiang2019fantastic, blanc2020implicit, liu2022pretraining, li2022happens, ma2022quadratic, lyu2023understanding, andriushchenko2023sgdlargestepsizes}. 

(3) Connecting loss landscape with feature learning: Some recent works~\citep{nakkiran2019sgd, rosenfeld2023outliers} tried to understand how the loss landscape is formed through the lens of feature learning. \citet{rosenfeld2023outliers} showed that a large learning rate will cause oscillation in learning subtle classification rules while continuing to learn other more deterministic features.~\citet{wang2024improvinggeneralizationconvergenceenhancing} studied how to improve generalization and convergence by amplifying the update provided by the optimizer in the flat direction of the loss landscape. ~\citet{wu2024largestepsizegradientdescent, cai2024largestepsizegradientdescent} studied gradient descent dynamics on logistic regression, showing that a large learning rate will cause oscillation in the earlier phase but will lead to higher progress later in training. \citet{pagliardini2024ademamixoptimizerbetterfaster} developed a modification of the Adam optimizer based on optimization analysis on the Rosenbrock function, which is a special case of the river valley landscape. \citet{song2024doessgdreallyhappen} shows that when SGD update is projected to the dominant subspace of the Hessian, the model's optimization progress slows down and they conjecture the existence of \emph{ill-conditioned valley} in the landscape, which can be viewed as a similar and simpler version of the river valley landscape discussed in this paper. 

(4) Ravines in the Loss Landscape. Concurrently with our work, \citet{davis2024gradientdescentadaptivestepsize} identified the existence of a ravine in the loss landscape—a manifold where every point has a vanishing gradient within the sharp eigenspace of the Hessian. This feature appears in any smooth loss function exhibiting fourth-order growth near minimizers. They also demonstrate the advantages of using adaptive step sizes in this context. The concept of a ravine aligns closely with the river structure described in our paper and can be considered a specific instance of it.

